\section{跨章节综合：如何选择自己的 architecture recipe}

前面逐项讨论了 norm、FFN、position、head shape、regularization 与 attention 变体，本节把这些局部选择压缩成一张决策表。核心不是给出永远正确的 recipe，而是区分稳定默认值、受推理系统驱动的变体，以及出现训练不稳定或长上下文需求时必须重新实验的边界。
\begin{longtable}{p{0.20\textwidth}p{0.32\textwidth}p{0.38\textwidth}}
\toprule
\textbf{决策层} & \textbf{保守默认值} & \textbf{什么时候需要重新实验} \\
\midrule
Norm & pre-norm + RMSNorm + no bias & 极深模型、训练不稳定、需要 QK norm 或额外 postnorm 时 \\
FFN & SwiGLU / GeGLU，按参数量调整 hidden dim & 小模型、特殊硬件、不同 activation kernel 性能差异明显时 \\
Position & RoPE & 超长上下文 extrapolation、局部 attention、特殊多模态输入时 \\
Head shape & 总 head dim 接近 model dim，head dim 取硬件友好值 & 推理 KV cache 成本高，考虑 GQA/MQA/MLA 时 \\
Regularization & 预训练少用 dropout，保留 weight decay & 数据很小、fine-tuning、过拟合明显时 \\
Stability & 监控 logits、QK norm、gradient norm & 出现 loss spikes、softmax 饱和、长上下文训练不稳时 \\
\bottomrule
\end{longtable}
\begin{importantbox}{最终 takeaways}
现代 LM 架构已经形成一组强默认值：pre-norm、RMSNorm、SwiGLU、RoPE、bias-free linear layers、保守 FFN/head/aspect ratio、GQA 或 hybrid attention 的推理友好变体。但这些默认值不是理论定律，而是公开模型、硬件、数据和训练稳定性共同塑造的经验均衡。
\end{importantbox}
\begin{warningbox}{不要把架构选择和系统选择分开}
RMSNorm、bias removal、parallel layers、GQA、sliding-window attention 看起来是 architecture，但常常由 memory bandwidth、kernel fusion、KV cache、latency 和 parallelism 驱动。Lecture 3 和 systems/inference 单元是相互嵌套的。
\end{warningbox}
